\begin{figure}
  \centering
  \begin{tikzpicture}[x=1.7cm,y=.62cm]
    \fill[hamblue!4] (0,0)--(4,4)--(4,0)--cycle;
    \draw[ham/arrow] (-.15,0)--(4.5,0) node[right] {insertions};
    \draw[ham/arrow] (0,-.25)--(0,4.85) node[above] {degree};
    \draw[dashed,hamgold,thick] (0,4)--(4.25,4);
    \node[anchor=west,font=\small] at (4.3,4) {$N=4$};
    \draw[dotted,thick] (0,0)--(1,1)--(2,2)--(3,3)--(4,4);
    \node[font=\small,fill=white] at (2.6,3.25) {$r(\ell)=\ell$};
    \draw[hamblue,line width=1.2pt]
    (0,0)--(1,1)--(2,2)--(3,1)--(4,2);
    \foreach \hamPointX/\hamPointY in {0/0,1/1,2/2,3/1,4/2} {
    \filldraw[draw=hamblue,fill=white] (\hamPointX,\hamPointY) circle (2pt);
    }
    \foreach \hamPointX in {1,2,3,4} {\node[font=\small,anchor=north] at (\hamPointX,-.18) {$\hamPointX$};}
    \node[ham/note,text width=8cm] at (2,-2.25)
    {Every prefix stays below its occurrence count.\\
    Projection changes none of these selected finite words.};
  \end{tikzpicture}
  \caption[Exact cutoffs for selected Gaussian words]{One possible degree path
    for four multiplications in a single Gaussian color. Multiplication is a sum
    of raising and lowering terms; the solid line illustrates one term, not the
    entire polynomial. For each color, every prefix has degree at most its
    occurrence count, which is at most the chosen cutoff. The result concerns
    selected finite words acting on the degree-zero vector, not entire
    propagators at a fixed cutoff. See \ProofRef{Lemma}{L:finite-word-cutoff}.}
  \label{Fig:hermite-cutoff}
\end{figure}
